% Regime map of rho = P_sinc/P_2s: kR versus cap half-angle for a smooth (v_0) and an
% edge-dominated (v_7) velocity pattern. Lightness-monotone map: light = sinc under-predicts,
% dark = sinc over-predicts. Cells are built in TeX from data/f8_*.dat.
\documentclass[border=1mm]{standalone}
\usepackage{amsmath,amssymb,bm}
\usepackage{jsvfigs}
\usepackage{pgfplots}
\pgfplotsset{compat=1.18}
\usepgfplotslibrary{patchplots}
\usepackage{pgfplotstable}

% lightness falls monotonically from rho = 1/8 (light) to rho = 8 (dark)
\definecolor{rhoa}{rgb}{0.88,0.93,0.99}
\definecolor{rhob}{rgb}{0.78,0.86,0.95}
\definecolor{rhoc}{rgb}{0.70,0.70,0.70}
\definecolor{rhod}{rgb}{0.55,0.30,0.28}
\definecolor{rhoe}{rgb}{0.15,0.04,0.04}

\pgfplotsset{
  colormap={rhomono}{color(-3)=(rhoa) color(-1.5)=(rhob) color(0)=(rhoc)
                     color(1.5)=(rhod) color(3)=(rhoe)},
  mapaxis/.style={
    scale only axis, width=49mm, height=30mm,
    xmin=0.09, xmax=42, ymin=-0.5, ymax=5.5, y dir=reverse,
    xtick={0.1,1,10,40}, xticklabels={0.1,1,10,40}, minor tick num=0,
    ytick={0,1,2,3,4,5}, xtick pos=bottom, ytick pos=left, tick align=outside, major tick length=1.2mm,
    axis line style={witness}, tick style={witness},
    every tick label/.append style={note},
    every axis label/.append style={lab},
    xlabel={$kR$},
    clip mode=individual,
    unbounded coords=jump,
    colormap name=rhomono, point meta min=-3, point meta max=3,
  },
}

\newcommand{\buildcells}[2]{%
  \pgfplotstablegetrowsof{#1}%
  \pgfmathtruncatemacro{\nr}{\pgfplotsretval-1}%
  \gdef#2{}%
  \foreach \i in {0,...,\nr}{%
    \pgfplotstablegetelem{\i}{x0}\of{#1}\let\xa\pgfplotsretval
    \pgfplotstablegetelem{\i}{x1}\of{#1}\let\xb\pgfplotsretval
    \pgfplotstablegetelem{\i}{row}\of{#1}\let\rw\pgfplotsretval
    \pgfplotstablegetelem{\i}{l2c}\of{#1}\let\mm\pgfplotsretval
    \pgfmathsetmacro{\ya}{\rw-0.5}\pgfmathsetmacro{\yb}{\rw+0.5}%
    \xdef#2{#2 (\xa,\ya) [\mm] (\xb,\ya) [\mm] (\xb,\yb) [\mm] (\xa,\yb) [\mm]}%
  }%
}

\begin{document}
\pgfplotstableread{data/f8_cells_cos0.dat}\tabA
\pgfplotstableread{data/f8_cells_cos7.dat}\tabB
\buildcells{\tabA}{\cellsA}
\buildcells{\tabB}{\cellsB}

\begin{tikzpicture}
% ---------------------------------------------------------------- (a) smooth pattern
\begin{axis}[mapaxis, xmode=log, name=axa, at={(0.9cm,2.45cm)}, anchor=south west,
  yticklabels={$15^\circ$,$30^\circ$,$60^\circ$,$90^\circ$,$120^\circ$,$150^\circ$}]
  \edef\tmp{\noexpand\addplot[patch, patch type=rectangle, shader=flat, point meta=explicit, draw=none]
     coordinates {\cellsA};}\tmp
  \addplot[draw=black, mark=none] table[x=x, y=y] {data/f8_rho1_cos0.dat};
  \node[lab, rotate=90, anchor=south] at ([xshift=-7.5mm]current axis.west) {cap half-angle $\theta_0$};
  % direct label for the rho=1 staircase
  \draw[leader] (axis cs:4.4721,5.5) -- ([xshift=-2.5mm,yshift=-2mm]axis cs:4.4721,5.5);
  \node[note, anchor=north east, inner sep=0.5pt, outer sep=0pt] at ([xshift=-2.5mm,yshift=-2mm]axis cs:4.4721,5.5) {$\rho=1$};
  % regime brackets (kR spans) and names
  \draw[witness] (axis cs:0.1,-0.55) -- (axis cs:0.1,-0.75) -- (axis cs:1,-0.75) -- (axis cs:1,-0.55);
  \draw[witness] (axis cs:1.5,-0.55) -- (axis cs:1.5,-0.75) -- (axis cs:8,-0.75) -- (axis cs:8,-0.55);
  \draw[witness] (axis cs:10,-0.55) -- (axis cs:10,-0.75) -- (axis cs:40,-0.75) -- (axis cs:40,-0.55);
  \node[note, anchor=south, align=center] at (axis cs:0.316,-0.85) {dipole regime,\\$\rho\gg1$};
  \node[note, anchor=south west] at (axis cs:1.5,-0.85) {minimum};
  \node[note, anchor=south, align=center] at (axis cs:20,-0.85) {two-face plateau,\\$\rho\approx0.5$};
  \node[paneltag, anchor=south west] at ([xshift=-1.2mm,yshift=11mm]current axis.north west) {(a)};
  \node[paneltitle, anchor=south west] at ([xshift=4mm,yshift=11mm]current axis.north west) {Smooth pattern $v_0$};
\end{axis}

% ---------------------------------------------------------------- (b) edge-dominated pattern
\begin{axis}[mapaxis, xmode=log, name=axb, at={(6.4cm,2.45cm)}, anchor=south west, yticklabels={}]
  \edef\tmp{\noexpand\addplot[patch, patch type=rectangle, shader=flat, point meta=explicit, draw=none]
     coordinates {\cellsB};}\tmp
  \addplot[draw=black, mark=none] table[x=x, y=y] {data/f8_rho1_cos7.dat};
  % edge-dominated corner: kR 10 to 40
  \draw[witness] (axis cs:10,-0.55) -- (axis cs:10,-0.75) -- (axis cs:40,-0.75) -- (axis cs:40,-0.55);
  \node[note, anchor=south east, align=right] at (axis cs:40,-0.85) {edge-dominated:\\sinc over-predicts};
  % near-closed cap, 150 degree row
  \draw[leader] (axis cs:42,5) -- ([xshift=3mm]axis cs:42,5);
  \node[note, anchor=west, text width=1.4cm] at ([xshift=3mm]axis cs:42,5)
     {near-closed cap: interior resonances};
  \node[paneltag, anchor=south west] at ([xshift=-1.2mm,yshift=11mm]current axis.north west) {(b)};
  \node[paneltitle, anchor=south west] at ([xshift=4mm,yshift=11mm]current axis.north west) {Edge-dominated pattern $v_7$};
\end{axis}

% ---------------------------------------------------------------- colour bar and key
\begin{scope}[shift={(4.2cm,1.05cm)}]
  \def\un{0.9}% width of one log2 step in cm; bar spans 6 steps
  \shade[left color=rhoa, right color=rhob] (0,0) rectangle (1.5*\un,0.25);
  \shade[left color=rhob, right color=rhoc] (1.5*\un,0) rectangle (3*\un,0.25);
  \shade[left color=rhoc, right color=rhod] (3*\un,0) rectangle (4.5*\un,0.25);
  \shade[left color=rhod, right color=rhoe] (4.5*\un,0) rectangle (6*\un,0.25);
  \draw[witness] (0,0) rectangle (6*\un,0.25);
  \foreach \t/\s in {-3/{1/8}, -2/{1/4}, -1/{1/2}, 0/{1}, 1/{2}, 2/{4}, 3/{8}} {
    \draw[witness] ({(\t+3)*\un},0) -- ({(\t+3)*\un},-0.08);
    \node[note, anchor=north, inner sep=0.5pt] at ({(\t+3)*\un},-0.08) {\s};
  }
  \node[note, anchor=east, align=right] at (-0.15,0.125) {sinc\\under-predicts};
  \node[note, anchor=west, align=left] at (6*\un+0.15,0.125) {sinc\\over-predicts};
  \node[note, anchor=north, align=center] at (3*\un,-0.50) {$\rho=P_{\rm sinc}/P_{2s}$, clipped at 1/8 and 8};
\end{scope}
\draw[black] (0.9,0.28) -- (1.5,0.28);
\node[note, anchor=west] at (1.55,0.28) {$\rho=1$ boundary};
\end{tikzpicture}
\end{document}


